\scanpage{210}
\mspage{64 bis}
impliquant les produits $\xi\times\eta$ ($\xi,\eta$ dans $A$) est satisfaite, alors pour $\forall$ flèche $u\colon I'\mathbin{\unsure{-|\to}}I$ dans $\mathrm{Diag}$, \unsure{la|$u$|elle} est $\mathbb{D}$-connexe ssi $u$ est $\mathbb{D}$-biasphérique.

\uline{Dém} L'\unsure{hypothèse}\deleted{\unsure{sur}} de conservativité, ou de fidélité, se transporte aux sommes indiciées par un ens. \deleted{\unsure{fini}} quelconque \inserted{(sous réserve d'existence)}, comme on voit \deleted{\unsure{par récur. sur le cardinal}.}\inserted{en écrivant $\coprod_\lambda\xi_\lambda=\xi_{\lambda_0}\amalg\coprod_{\lambda\ne\lambda_0}\xi_\lambda$}. L'asser\deleted{\unsure{fide}}tion faite sur $u_!$ résulte \inserted{alors} du calcul de $u_!$ fibre par \unsure{fibres|fibre} (qui est valable vu que $u$ est étale donc lisse), en tenant compte que les fibres \unsure{s.x|sont} discrètes.

Je ne connais pas d'exemple d'un couple $(\mathbb{D},u)$ d'un dérivateur $\mathbb{D}$ et d'une flèche $u$ de $\mathrm{Diag}$, pour lesquels $u$ soit $\mathbb{D}$-connexe sans être $\mathbb{D}$-asphérique. On a donné des conditions sur $\mathbb{D}$, \deleted{\unsure{ou sur $I=$ but}}, qui assurent l'équivalence des deux notions. Si $\mathbb{D}$ est de la forme $\mathbb{D}_M$, et si \deleted{\unsure{donc}} $M$ satisfait la condition dite dans \unsure{4.7|4.5|4.?} ou sa duale, alors par \unsure{4.7|4.?} c'est O.K. Donc un contrexemple avec un $\mathbb{D}_M$ devrait en tenir compte. Je ne vais pas \unsure{m'obstiner} pour tirer au clair ce qui en est.
% archive sheet number: 210
